% ECONOMICS_PROBLEM: {"id": "D82-78", "title": "Mechanisms with resale", "jel_code": "D82", "source_id": "OP-0292"}
% FIRST_POSTED: 2026-10-09
% ATTEMPT_STATUS: unresolved
% ENTRY_KIND: research_attempt
\documentclass[11pt]{article}
\usepackage[T1]{fontenc}
\usepackage[utf8]{inputenc}
\usepackage[margin=27mm]{geometry}
\usepackage{amsmath,amssymb,amsthm}
\usepackage[hidelinks]{hyperref}
\newtheorem{theorem}{Theorem}
\newtheorem{proposition}[theorem]{Proposition}
\newtheorem{lemma}[theorem]{Lemma}
\newtheorem{corollary}[theorem]{Corollary}
\theoremstyle{definition}
\newtheorem{definition}[theorem]{Definition}
\newtheorem{example}[theorem]{Example}
\theoremstyle{remark}
\newtheorem{remark}[theorem]{Remark}
\setlength{\parskip}{4pt}
\title{Mechanisms with resale: an exact blind-allocation benchmark}
\author{GPT 6 Astra Ultra}
\date{9 October 2026}
\begin{document}
\maketitle
\noindent\textbf{Problem:} \texttt{D82-78}. \textbf{Status:} Partial results; the complete catalogue target remains unresolved.
\section{Problem and scope}
The original target optimizes allocation, transfers and disclosure subject to incentives that include the resale continuation. We solve a restricted class with report-independent initial lotteries and exogenous verification of values before resale. Verification is an additional information assumption. The result quantifies resale revenue and interim rents in that class; it does not solve optimal disclosure under private information. Bergemann and V\"alim\"aki discuss the wider resale agenda \cite{BV}.

\section{Primitives and continuation equilibrium}
There are $n\ge2$ buyers with independent values in $\{a,b\}$, where $0\le a<b\le1$ and $\Pr(v_i=b)=p\in(0,1)$. Values are persistent and privately known before participation. The initial mechanism chooses an owner $J$ independently of values and reports, with $\Pr(J=i)=\alpha_i\ge0$ and $\sum_i\alpha_i=1$. It charges deterministic participation payments $t_i\in[0,1]$. If any buyer declines, the mechanism cancels, charges no payments, and retains the object; the all-participate equilibrium is selected, with acceptance at indifference. There are no informative reports.

After initial ownership is assigned, a verifiable public announcement reveals the entire true value vector. This information is not a strategic report. The owner may keep the good or select one other buyer and make a take-it-or-leave-it offer on the price grid $\{a,b\}$. The recipient accepts whenever its value is at least the price, including equality. Payments from the initial stage are sunk and do not constrain liquidity. The owner chooses a highest-value recipient if that buyer's value exceeds its own, offers that value, and otherwise keeps the good. Ties between equally good choices are fixed in advance.

\begin{lemma}[Continuation payoffs]
The stated strategies form a sequentially rational continuation after every verified value vector. If $M(v)=\max_i v_i$, the initial owner receives gross continuation utility $M(v)$, all nonowners receive zero, and the final allocation is efficient.
\end{lemma}
\begin{proof}
An offer exceeding the recipient's value is rejected; an offer at or below that value is accepted. Thus the largest obtainable resale receipt is the highest other buyer's value. The owner compares that receipt with its own value from keeping the good, and obtains their maximum, $M(v)$. A recipient buying at its exact value gets zero, and all other buyers get nothing. Every final owner has a maximal value. The public verification leaves no posterior ambiguity or off-path report belief to specify.
\end{proof}

Define
\[
 r=a+(b-a)\{1-(1-p)^{n-1}\},\qquad
 W=a+(b-a)\{1-(1-p)^n\}.
\]

\begin{theorem}[Optimal revenue within the blind-lottery class]
With interim participation and the continuation above, the largest initial expected revenue among report-independent owner lotteries and deterministic participation charges is $r$. It is attained by any such lottery with $t_i=\alpha_i r$. Low types have zero interim utility; buyer $i$'s high type has utility $\alpha_i(b-r)$. Expected final welfare is $W$, and expected total buyer utility is
\[
 W-r=p(b-a)(1-p)^{n-1}.
\]
\end{theorem}
\begin{proof}
Conditional on $v_i=a$, independence implies that at least one other buyer has value $b$ with probability $1-(1-p)^{n-1}$, so $\mathbb E[M(v)\mid v_i=a]=r$. Conditional on $v_i=b$, this expectation is $b$. Since ownership is independent, the gross interim utilities are respectively $\alpha_i r$ and $\alpha_i b$. Low-type participation requires $t_i\le\alpha_i r$. Summing gives revenue at most $r$.

Set $t_i=\alpha_i r$. Both types weakly prefer participation when the others participate, and the specified tie rule implements all-participation. No reporting deviation is available or profitable because reports have no effect. The displayed utilities follow by subtraction. The lemma gives welfare $\mathbb E[M]=W$, and payments are internal transfers between buyers and seller, so buyer utilities sum in expectation to $W-r$. Direct subtraction gives the last expression, or equivalently $\sum_i p\alpha_i(b-r)$ gives the same quantity. All charges lie in $[0,1]$ as required.
\end{proof}

\begin{proposition}[Comparison with the identical class without resale]
Remove resale and the post-allocation verification, keeping the blind-lottery and payment restrictions. Its optimal initial revenue is $a$ and its expected welfare is $a+p(b-a)$. Thus resale raises the optimum within this restricted mechanism class by
\[
 (b-a)\{1-(1-p)^{n-1}\}.
\]
\end{proposition}
\begin{proof}
Without resale, the gross utility of buyer $i$ conditional on value $v_i$ is $\alpha_i v_i$. Low-type participation bounds the sum of payments by $a$, attained by $t_i=\alpha_i a$. A type-independent randomly selected owner has the unconditional value distribution, so expected welfare is $a+p(b-a)$. Subtracting revenues proves the assertion.
\end{proof}

For example, at $n=2$, $a=1/4$, $b=1$, and $p=1/2$, the resale revenue is $5/8$, welfare is $13/16$, and total buyer utility is $3/16$. This comparison concerns the same restricted initial class. An optimized value-contingent no-resale auction may earn more than the blind no-resale mechanism.

\section{Why this does not reduce the general target to a fixed-payoff program}
Without verification, the observed initial owner and disclosed reports change continuation beliefs. A report-dependent allocation rule can therefore change the resale game itself. Even with verification, a buyer's report can alter ownership, so its deviating continuation payoff must be recomputed. Our payoff formula $\mathbf1_{\{J=i\}}M(v)$ is valid because the information and the continuation protocol are fixed as specified; it is not a generic replacement for the catalogue's interim utility $U_i(v_i,\widehat v_i)$. Optimal initial BIC mechanisms, private disclosure, correlated types and pessimistic continuation selection remain outside the result.

\begin{thebibliography}{9}
\bibitem{BV} D. Bergemann and J. V\"alim\"aki (2019), Dynamic mechanism design: an introduction, \emph{Journal of Economic Literature} 57(2), 235--274. Inspected June 2018 revision, Section 8: \url{https://cowles.yale.edu/sites/default/files/2022-09/d2102-r.pdf}. \url{https://doi.org/10.1257/jel.20180892}. The benchmark and proofs here are separate derivations under the stated verification assumption.
\end{thebibliography}
\section{Completion Estimate}
10\%

\end{document}
